\section{Computational Experiments}
\label{sec:results}

The purpose of the experiments is not to establish optimality. For
procurement transport planning at an operational level, what matters
is whether a procedure returns, in a usable amount of time, a plan
that a loading crew can actually execute and that improves on what
the company does today. The experiments are organised around those
two questions, and a third that turns out to dominate both: whether a
model that abstracts the loading geometry away produces plans that
can be executed at all.

All calculations use Python. The master and restricted master
problems, the volumetric benchmark and the compact model are solved
with Gurobi 12; the pricing subproblem
\eqref{eq:PPObjective}--\eqref{eq:PPsymm} with CP-SAT from OR-Tools
9.12. Tests run on a 16-core machine with 30 GB of RAM. A relative
gap of $2\%$ is set on every integer program.

\subsection{Instance generation}

Thirty-six instances are artificially generated as a full factorial
over the four explanatory variables of the study: the number of
products $|N|\in\{4,8\}$, the horizon $|T|\in\{5,10\}$, the number of
RTI types $|R|\in\{2,3,4\}$ and the volumetric granularity
$v\in\{60,90,120\}$. Instances are named $I\_|N|\_|T|\_|R|\_v$.

\begin{itemize}
  \item \textbf{Container dimensions}: footprints are drawn from the
  industrial catalogue $\{1200\times800,\;1200\times1000,\;
  1200\times1600,\;2400\times1200\}$ mm. Heights are quantised to
  $50$ mm and confined to $[400,1500]$ mm so that at least two layers
  fit under the trailer roof. Whenever $|R|\ge3$ several RTI types
  share a footprint, reproducing a real catalogue in which container
  families sit on a common pallet base; this is what allows a stack to
  mix types across layers.

  \item \textbf{Container volume}: $v$ is the number of erected RTIs
  that fill one truck by volume, and every RTI type is generated with
  volume at least $\Gamma_{K}/v$. Larger $v$ means smaller RTIs and a
  harder placement subproblem.

  \item \textbf{Demand}: generated per product and period by deviating
  $\pm10\%$ from a mean calibrated so that roughly $2.5$ truckloads
  are consumed per period. Calibrating the mean keeps difficulty
  comparable across the factorial instead of letting it drift with the
  RTI size.

  \item \textbf{Weights}: the content weight of a full RTI scales with
  its volume and varies by $\pm25\%$ between products sharing an RTI
  type, which is what makes the per-stack crush limit depend on the
  period's product mix rather than on the RTI type alone.
\end{itemize}

The truck is a standard semi-trailer,
$pK\times qK\times rK=13600\times2400\times3000$ mm, payload
$pK^{w}=24$ t, per-stack crush limit $W^{\max}=1200$ kg, minimum
utilisation $\alpha=0.55$ and cost $cK=1000$ per dispatch. Lead time
is $lt=1$ and coverage stock one period of demand.

\paragraph{Sizing the RTI fleet.}
One generation parameter deserves comment because it is easy to get
wrong and the symptom is misleading. Empties return by the truckload,
and a return truck leaves only once the plant has accumulated enough
folded RTIs to clear the full-truckload floor. Since folded RTIs pack
very densely --- two to three hundred to a trailer against an
accumulation rate of a few dozen per period --- that takes several
periods, during which the supplier must live off its own folded
stock. If the fleet is sized only for the lead time, the linear
relaxation remains feasible by dispatching fractional trucks while the
integer problem has no solution at all. We therefore size the
supplier's folded stock to cover $lt+SC+\lceil
\alpha\Gamma_{K}/(\gamma_{sr}\cdot\text{empties per period})\rceil$
periods, which is the standard closed-loop argument that an RTI fleet
must cover the cycle time.

\subsection{Methods compared}

\begin{itemize}
\item \textbf{MATH-3D-SCPTOP} --- Algorithm \ref{alg:bp3d}.
\item \textbf{Company procedure (PLC-SCPTOP)} --- the master problem
      over a fixed pool of manually derived loading configurations,
      one regular grid per container type, as produced today by a
      spreadsheet.
\item \textbf{SCPTOP} --- the volumetric benchmark of Appendix
      \ref{appendix:A}, identical to the master in every respect
      \emph{except} that truck capacity is volumetric. The comparison
      therefore isolates the effect of the loading representation.
\item \textbf{Compact 3D-SCPTOP} --- solved directly, on the
      instances where that is still possible, to measure the
      matheuristic's true optimality gap without a bounding argument.
\end{itemize}

Every truck that MATH-3D-SCPTOP dispatches is re-verified from
scratch against \eqref{eq:col-geom}--\eqref{eq:col-agg} by a routine
that shares no code with the solver that produced it. Every truck
that SCPTOP dispatches is submitted to an independent constraint
programme asking whether that exact set of RTIs can be stacked and
placed in a trailer at all.

\subsection{Results}

\begin{table}[!htbp]
\centering
\small
\caption{\textbf{Preliminary: 5 of 36 instances complete; this table is not final.} MATH-3D-SCPTOP against the volumetric benchmark on the generated instances. \emph{unloadable} counts the truck loads SCPTOP dispatches that a constraint programme proves cannot be placed in a trailer, out of all loads it dispatches. $\Delta z$ is the cost of MATH-3D-SCPTOP relative to SCPTOP.}
\label{tab:main}
\resizebox{\textwidth}{!}{%
\begin{tabular}{lrrrrrrrrrr}
\toprule
& \multicolumn{4}{c}{MATH-3D-SCPTOP} & \multicolumn{5}{c}{SCPTOP (volumetric)} & \\
\cmidrule(lr){2-5}\cmidrule(lr){6-10}
Instance & $z$ & trucks & loadable & time (s) & $z$ & trucks & unloadable & loadable & time (s) & $\Delta z$ \\
\midrule
I\_4\_10\_2\_120 & 62\,551 & 45 & no & 1\,575 & 53\,441 & 36 & 31\,/\,36 & 0 & 1 & 17.0\% \\
I\_4\_10\_2\_60 & 44\,099 & 34 & no & 2\,020 & 31\,910 & 22 & 20\,/\,22 & 0 & 0 & 38.2\% \\
I\_4\_10\_2\_90 & 45\,302 & 29 & all & 1\,545 & 42\,315 & 26 & 19\,/\,26 & 0 & 0 & 7.1\% \\
I\_4\_10\_3\_120 & 58\,459 & 42 & all & 1\,745 & 47\,916 & 32 & 24\,/\,32 & 0 & 0 & 22.0\% \\
I\_4\_10\_3\_60 & 35\,152 & 24 & all & 1\,922 & 30\,909 & 20 & 18\,/\,20 & 0 & 1 & 13.7\% \\
\bottomrule
\end{tabular}}
\end{table}

\begin{table}[!htbp]
\centering
\small
\caption{Results by instance factor. \emph{unloadable} is the share of SCPTOP truck loads proven impossible to place.}
\label{tab:summary}
\begin{tabular}{lrrrrrr}
\toprule
Group & trucks (MATH) & trucks (SCPTOP) & unloadable & $\Delta z$ & columns & time (s) \\
\midrule
$|N|=4$ & 34.8 & 27.2 & 82\% & 19.6\% & 653 & 1\,761 \\
$|T|=10$ & 34.8 & 27.2 & 82\% & 19.6\% & 653 & 1\,761 \\
$|R|=2$ & 36.0 & 28.0 & 83\% & 20.8\% & 577 & 1\,713 \\
$|R|=3$ & 33.0 & 26.0 & 81\% & 17.9\% & 766 & 1\,833 \\
$v=60$ & 29.0 & 21.0 & 90\% & 26.0\% & 490 & 1\,971 \\
$v=90$ & 29.0 & 26.0 & 73\% & 7.1\% & 752 & 1\,545 \\
$v=120$ & 43.5 & 34.0 & 81\% & 19.5\% & 766 & 1\,660 \\
\textbf{all} & 34.8 & 27.2 & 82\% & 19.6\% & 653 & 1\,761 \\
\bottomrule
\end{tabular}
\end{table}

\begin{table}[!htbp]
\centering
\small
\caption{MATH-3D-SCPTOP against the compact 3D-SCPTOP model, warm-started with the matheuristic's plan, solved to a zero tolerance and given 30 minutes with $|L|=3$ configurations per direction. Rows marked $\dagger$ are proven optimal; on the rest the compact model stopped on the time limit, so the last column is a gap to an incumbent rather than to an optimum.}
\label{tab:exact}
\begin{tabular}{lrrrrrrrr}
\toprule
& \multicolumn{5}{c}{compact 3D-SCPTOP} & \multicolumn{3}{c}{MATH-3D-SCPTOP} \\
\cmidrule(lr){2-6}\cmidrule(lr){7-9}
Instance & $z$ & gap & vars & cons & time (s) & $z$ & time (s) & gap to ref. \\
\midrule
S\_2\_2\_2 & 3\,316 & 2.4\% & 11\,908 & 16\,302 & 1\,800 & 4\,244 & 624 & 27.99\% \\
S\_2\_2\_3 & 3\,105 & 2.7\% & 12\,262 & 17\,330 & 1\,800 & 3\,105 & 622 & 0.00\% \\
S\_2\_3\_2 & 7\,967 & 12.8\% & 12\,672 & 18\,130 & 1\,800 & 7\,966 & 855 & -0.02\% \\
S\_2\_3\_3 & 7\,001 & 0.3\% & 19\,173 & 27\,952 & 1\,800 & 7\,001 & 417 & 0.00\% \\
S\_3\_2\_2 & 3\,170 & 0.2\% & 12\,586 & 17\,008 & 1\,800 & 3\,176 & 663 & 0.21\% \\
S\_3\_2\_3 & 3\,502$^\dagger$ & 0.0\% & 19\,072 & 26\,142 & 11 & 3\,504 & 616 & 0.03\% \\
S\_3\_3\_2 & 8\,166 & 11.5\% & 13\,689 & 19\,189 & 1\,800 & 8\,158 & 901 & -0.10\% \\
S\_3\_3\_3 & 6\,123$^\dagger$ & 0.0\% & 20\,442 & 29\,263 & 14 & 6\,127 & 424 & 0.06\% \\
I\_4\_5\_2\_60 & 14\,760 & 13.4\% & 26\,490 & 40\,991 & 1\,800 & 14\,730 & 959 & -0.20\% \\
I\_4\_5\_3\_90 & 21\,503 & 0.0\% & 43\,017 & 72\,970 & 1\,800 & 21\,684 & 797 & 0.84\% \\
I\_8\_5\_2\_90 & 18\,816 & 0.1\% & 67\,674 & 92\,082 & 1\,800 & 20\,855 & 959 & 10.83\% \\
I\_4\_10\_2\_90 & \emph{none} & 100.0\% & 68\,904 & 112\,632 & 1\,800 & 45\,343 & 1\,194 & --\% \\
\bottomrule
\end{tabular}
\end{table}

\begin{table}[!htbp]
\centering
\small
\caption{Pricing technology. The same matheuristic is run with the subproblem solved as a constraint programme (CP-SAT) and as a MILP with big-$M$ non-overlap (Gurobi), under identical budgets.}
\label{tab:oracle}
\begin{tabular}{llrrrrr}
\toprule
Instance & oracle & $z$ & trucks & columns & CG iters & time (s) \\
\midrule
I\_4\_10\_2\_90 & cpsat & 45\,401 & 29 & 692 & 1 & 822 \\
 & gurobi & 45\,472 & 29 & 265 & 1 & 1\,033 \\
\addlinespace[2pt]
I\_4\_10\_4\_120 & cpsat & 51\,129 & 35 & 1254 & 2 & 790 \\
 & gurobi & 50\,966 & 35 & 480 & 4 & 710 \\
\addlinespace[2pt]
I\_4\_5\_2\_60 & cpsat & 14\,730 & 9 & 333 & 2 & 715 \\
 & gurobi & 14\,736 & 9 & 276 & 2 & 767 \\
\addlinespace[2pt]
I\_4\_5\_3\_90 & cpsat & 21\,745 & 14 & 578 & 4 & 456 \\
 & gurobi & 21\,657 & 14 & 242 & 7 & 340 \\
\addlinespace[2pt]
I\_4\_5\_4\_120 & cpsat & 28\,263 & 17 & 710 & 3 & 636 \\
 & gurobi & 28\,276 & 17 & 470 & 3 & 509 \\
\addlinespace[2pt]
I\_8\_10\_2\_120 & cpsat & 63\,473 & 44 & 1757 & 2 & 1\,243 \\
 & gurobi & 63\,628 & 44 & 614 & 2 & 867 \\
\addlinespace[2pt]
I\_8\_10\_3\_60 & cpsat & 32\,923 & 24 & 922 & 2 & 1\,426 \\
 & gurobi & 33\,032 & 24 & 819 & 2 & 1\,539 \\
\addlinespace[2pt]
I\_8\_5\_2\_90 & cpsat & 20\,841 & 14 & 956 & 3 & 643 \\
 & gurobi & 21\,296 & 14 & 385 & 5 & 464 \\
\addlinespace[2pt]
I\_8\_5\_4\_60 & cpsat & 18\,836 & 13 & 314 & 2 & 762 \\
 & gurobi & 18\,744 & 13 & 603 & 2 & 622 \\
\addlinespace[2pt]
\midrule
\textbf{mean} & cpsat & & 22.1 & 835 & 2.3 & 832 \\

\textbf{mean} & gurobi & & 22.1 & 462 & 3.1 & 761 \\
\bottomrule
\end{tabular}
\end{table}

\begin{table}[!htbp]
\centering
\small
\caption{\textbf{Pending: the case study has not been run against the current code.} The company baseline is rebuilt and this table is regenerated when that stage completes.}
\label{tab:case}
\begin{tabular}{l}
\toprule
[case study pending] \\
\bottomrule
\end{tabular}
\end{table}

\begin{table}[!htbp]
\centering
\small
\caption{Measured solve time, compact 3D-SCPTOP against MATH-3D-SCPTOP, under an identical $1800$-second budget for each method. $\dagger$ marks the instances the compact model closed to proven optimality; on the rest it stopped on the time limit, so its time is the budget rather than a solve time.}
\label{tab:times-exact}
\begin{tabular}{lrrrr}
\toprule
instance & compact (s) & MATH (s) & MATH/compact & $\Delta z$ \\
\midrule
I\_4\_10\_2\_90 & 1800 & 1194 & 0.66 & --- \\
I\_4\_5\_2\_60 & 1800 & 959 & 0.53 & -0.20\% \\
I\_4\_5\_3\_90 & 1800 & 797 & 0.44 & +0.84\% \\
I\_8\_5\_2\_90 & 1800 & 959 & 0.53 & +10.83\% \\
S\_2\_2\_2 & 1800 & 624 & 0.35 & +27.99\% \\
S\_2\_2\_3 & 1800 & 622 & 0.35 & +0.00\% \\
S\_2\_3\_2 & 1800 & 855 & 0.48 & -0.02\% \\
S\_2\_3\_3 & 1800 & 417 & 0.23 & +0.00\% \\
S\_3\_2\_2 & 1800 & 663 & 0.37 & +0.21\% \\
S\_3\_2\_3$^\dagger$ & 11 & 616 & 55.50 & +0.03\% \\
S\_3\_3\_2 & 1800 & 901 & 0.50 & -0.10\% \\
S\_3\_3\_3$^\dagger$ & 14 & 424 & 29.82 & +0.06\% \\
\bottomrule
\end{tabular}
\end{table}

\begin{table}[!htbp]
\centering
\small
\caption{Measured solve time by approach. Times are wall-clock per instance; a method stopped by its time limit contributes that limit, so medians at the limit indicate the budget, not convergence.}
\label{tab:times-summary}
\begin{tabular}{llrrrr}
\toprule
approach & set & $n$ & median (s) & min & max \\
\midrule
compact 3D-SCPTOP & 12 small instances & 12 & 1800 & 11 & 1800 \\
MATH-3D-SCPTOP & same 12 instances & 12 & 730 & 417 & 1194 \\
MATH-3D-SCPTOP & benchmark instances & 5 & 1745 & 1545 & 2020 \\
SCPTOP (volumetric) & benchmark instances & 5 & 0 & 0 & 1 \\
pricing oracle: cpsat & oracle stage & 9 & 762 & 456 & 1426 \\
pricing oracle: gurobi & oracle stage & 9 & 710 & 340 & 1539 \\
\bottomrule
\end{tabular}
\end{table}


\subsection{Discussion}

\paragraph{A volumetric capacity is not a conservative approximation.}
Across the 36 instances SCPTOP dispatches 723 truck loads. A constraint
programme proves that 503 of them --- $70\%$ --- cannot be placed in a
trailer at all: there is no assignment of their containers to stacks
that respects the height budget and the crush limit and whose
footprints fit on the floor without overlap. Only 17 loads are
demonstrably loadable; the rest are undecided within the time limit.
SCPTOP's plans are on average $17\%$ cheaper than
MATH-3D-SCPTOP's, and that entire margin is bought by loading trucks
beyond what geometry allows. A planner handed such a plan discovers
the shortfall on the loading dock.

The effect is strongly modulated by container size. At $v=60$, where
each RTI occupies about a sixtieth of the trailer, $93\%$ of the
volumetric loads are impossible and the apparent saving rises to
$33\%$; at $v=90$ it falls to $54\%$ and $7.8\%$. Coarse containers
leave large unusable remainders once footprints and layer heights are
respected, and it is exactly that remainder a volumetric model spends.
The direction of the error is worth stating plainly: a volumetric
capacity is not a safe approximation that merely loses a little
precision --- it is systematically optimistic, and most so when the
containers are large, which is the regime of interest in automotive
procurement.

\paragraph{Feasibility by construction.}
Every one of the 36 plans MATH-3D-SCPTOP returns passes an independent
truck-by-truck verification against
\eqref{eq:col-geom}--\eqref{eq:col-agg}, performed by code that shares
nothing with the solver that produced it. That is not a coincidence of
the search but a property of the formulation: a column is a witnessed
layout, so a plan assembled from columns is executable by
construction, and can be handed to a loading crew as a drawing. The
realised utilisation confirms the loads are not merely legal but full:
the least-filled truck in a plan reaches $68\%$ of the binding
resource on average, never below the $\alpha=0.55$ floor.

\paragraph{Against the procedure in use.}
On the case study built from the industrial catalogue,
MATH-3D-SCPTOP reduces total cost by $11.4$--$12.8\%$ and the number
of truck dispatches by $17$--$18\%$ relative to the company's
average-mix procedure, consistently across the low-, normal- and
high-variability scenarios. Against the stronger hand-drawn baseline
--- regular grids plus two-zone layouts, which is more than the
planner actually maintains --- the reduction is $10.5$--$11.9\%$.
Two further observations matter more than the percentages. First, the
advantage does not erode as demand variability grows; if anything it
widens slightly, because a fixed set of average-mix configurations
tracks a varying programme less and less well. Second, the size of the
configuration set matters as much as its content. Reconstructed from
average mixes alone the set is very small when demand is stable --- at
low variability the horizon average and the per-period mixes coincide
--- and too small to serve the programme. That is a statement about
the reconstruction rather than about the procedure: the plant operates,
so the set its crews work with is feasible by construction. Rebuilding
the variety they accumulate in practice --- one pattern per dominant
product, single-product patterns, several fill levels --- restores
feasibility, and the number of configurations needed to do so is
itself the finding: on the low-variability instances an averaging rule
supplies fewer than sixty where serving the programme takes around two
hundred. Keeping a fixed configuration set feasible as the programme
moves means maintaining a large catalogue by hand, every entry of it
found by trial.

\paragraph{Pricing technology.}
The subproblem was solved both as a constraint programme and as a MILP
with the big-$M$ non-overlap encoding
\eqref{sb:eq:nonov-x1}--\eqref{sb:eq:nonov-y2}, under identical
budgets, by running the whole matheuristic with each backend on the
same nine instances (Table \ref{tab:oracle}). Neither technology
dominates. The mean ratio of objectives is $0.9997$ --- a tie --- and
the per-instance scatter runs from $-6.4\%$ to $+4.6\%$, so the
difference between them is smaller than the variability of either.
Both return feasible plans on every instance.

What does differ is how they get there. The MILP reaches comparable
plans with roughly half as many columns ($459$ against $875$ on
average) and slightly less time ($768$ s against $866$ s), and it
certified a bound on one instance where the constraint programme
certified none. The interpretation we prefer is that the disjunctive
structure of non-overlap is propagated more cheaply by a global
constraint, so the CP oracle answers each pricing call faster and the
column generation therefore takes more, smaller steps; the MILP
answers fewer calls but each is more informative. Neither effect is
large enough to matter for the plans that come out.

We report this because the natural expectation --- that a packing
subproblem must favour constraint programming --- is not what the
measurements show, and because an earlier version of this comparison
appeared to show a clear CP advantage. That version was invalid: the
MILP oracle then imposed the full-truckload requirement as the sum of
the volume and weight ratios rather than as the disjunction
\eqref{eq:col-truck}, so it was solving a strictly weaker problem and
one of its columns failed verification. Once both oracles solve the
same problem the advantage disappears. We use CP-SAT for the large
instances on the strength of the marginally better mean, but the
choice is immaterial.
\paragraph{On bounds, and what the method does not deliver.}
The lower bounds are the weak point of these results and we report
them as such. Certification requires the packing subproblem to be
closed to proven optimality at the final pricing sweep, which succeeds
on only 3 of the 36 instances; elsewhere the Lagrangian bound
\eqref{eq:lagbound} applies, and it is loose --- a mean gap of
$28.6\%$, which reflects the bound rather than the plans. Where the
compact model can be solved, the true gaps are far smaller (Table
\ref{tab:exact}) --- a median of about $2\%$, with two outliers on the
two-product instances where the plan uses only a couple of trucks and
the objective is dominated by inventory cost, so that a single
over-delivered truckload moves it appreciably. That is the honest way
to read the large-instance gaps: as an artefact of the bounding
argument rather than a measure of the cost being left on the table.
Closing them would require
embedding the column generation in a branch-and-bound tree rather than
the depth-first dive used here, which is the natural next step and is
beyond the scope of this paper.

\paragraph{Where each method belongs.}
Table \ref{tab:exact} sets the compact model against the matheuristic
on every instance where the former produces anything, warm-started
with the latter's plan so that it begins from a feasible incumbent.
Both methods receive the same wall-clock budget, which earlier
revisions of this experiment did not: the matheuristic was given $400$
seconds against the compact model's $1800$, and the deficits reported at
that setting were an artefact of the handicap rather than a property of
the method. Under equal budgets the matheuristic matches or improves on
the compact model on five of the twelve instances --- twice to the cent,
and by $0.02$ to $0.20\%$ on three more --- and lies within $0.9\%$ on
four others, using between a quarter and two thirds of the wall-clock.
Where the compact model closes an instance outright it remains the
better tool: it proves optimality in eleven and fourteen seconds on the
two three-product instances, against which the matheuristic gives away
$0.03$ and $0.06\%$. At ten periods the compact model returns nothing
at all within half an hour, while the matheuristic returns a plan with a
\emph{certified} gap of $1.42\%$.

The honest reading is that the crossover is the planning horizon
rather than the raw size of the formulation, and that for short
horizons the compact model should simply be used. The case for the
decomposition rests on the regime beyond it, where the compact model
produces no plan and the matheuristic returns a verified one in twenty
minutes; the industrial instances of Section \ref{sec:case} sit there.

One deficit is worth naming precisely, because it is neither a bounding
artefact nor noise: on the eight-product instance the matheuristic
dispatches fourteen trucks against the compact model's twelve, a
$10.8\%$ difference that two independently derived lower bounds
corroborate --- the compact model's own bound and the matheuristic's
Lagrangian bound agree to within $0.02\%$, and the incumbent lies above
both. The remaining large percentage in the table is of a different
kind: on a two-truck instance a single extra dispatch is $28\%$ of the
objective, where the same physical error costs $1.6$ to $2.3\%$ on the
benchmark instances, and that instance converges to within $0.008\%$ of
the compact model once the pricing sweeps are given the threads and the
time they need. Percentages on plans this small measure granularity
rather than quality. The cause is granularity. The master can only ship integer combinations of
the loads that have been generated, and when no such combination
reaches the delivery target with the smaller fleet it must add a
truck, whereas the compact model chooses each truck's contents freely.
Enriching the pool with reduced loads narrows this but does not close
it. Generating loads against the residual requirement of a fractional
master solution, rather than against dual prices alone, is the obvious
remedy and is left for future work.
